\section{Background and Related Work}
\label{related}

In this section of the paper, we first provide a background on the definition of vulnerability and the different steps in software vulnerability detection. Then we discuss the related surveys and highlight their differences compared to our survey.

\subsection{Background}

Software vulnerability management is now essential for guaranteeing the security and integrity of software systems~\cite{chang2011trend, foreman2019vulnerability, raza2022threat, walkowski2021vulnerability}. Given the increasing reliance on software for many critical processes such as financial transactions, the frequency of vulnerabilities poses serious risks..~\cite{lu2021neucheck, he2020smart}, autonomous driving~\cite{gao2021autonomous, luo2019localization}, and mission-critical systems~\cite{goseva2017experience, harzevili2022characterizing}. Software vulnerabilities can be exploited by malicious entities to gain unauthorized access, compromise sensitive information, or disrupt services if they go undetected or ignored.~\cite{harzevili2022characterizing}. As a result, excellent software vulnerability management is crucial to handling these risks, preserving user privacy~\cite{anthonysamy2017privacy}, maintaining system availability, and assuring software application trustworthiness~\cite{medeiros2023trustworthiness}. By proactively detecting, analyzing, and remediating vulnerabilities, organizations may strengthen their software systems against changing cybersecurity threats~\cite{aslan2023comprehensive} and adhere to industry best practices for safe software development and deployment.


There are multiple steps in software vulnerability management including vulnerability detection~\cite{cao2022mvd}, vulnerability analysis~\cite{kudjo2020effect}, and vulnerability remediation~\cite{le2021survey}. In the following subsections, we elaborate on each step in detail.

\subsubsection{Vulnerability detection}
Vulnerability detection is critical in the overall process of managing software vulnerabilities~\cite{cao2022mvd, jeon2021autovas,
 tian2020bvdetector,
 liu2020cd,
 wang2020combining,
 cheng2021deepwukong,
 zou2019mu,
 cao2022mvd,
 ziems2021security,
 lin2019software,
 li2021sysevr}. It comprises detecting and investigating possible security weaknesses in software systems that attackers may exploit. There are several traditional techniques commonly used for vulnerability detection:

\noindent \textbf{Manual Code Auditing} \cite{staron2020using, bacchelli2013expectations, carlsson2005software, shar2012auditing, shar2010auditing}: In this method, human experts examine the source thoroughly with the goal of manually detecting coding flaws, unsafe procedures, and possible vulnerabilities. Manual code review is time-consuming and requires the knowledge of qualified developers or security analysts. However, it provides for a thorough grasp of the code and can reveal subtle bugs that automated tools may overlook.

\noindent \textbf{Static Analysis}~\cite{infer, flawfinder, rats, cppcheck, SpotBugs, lattner2008llvm}: Static analysis involves using automated tools to analyze the source code or compiled binaries without executing the software. It examines the code structure, identifies potential coding issues, and detects common vulnerabilities such as buffer overflows~\cite{harzevili2022characterizing}, injection attacks, and insecure data handling. Static analysis tools employ various techniques like data flow analysis, control flow analysis, and pattern matching to identify potential vulnerabilities. They can help scale vulnerability detection efforts by analyzing large codebases efficiently.

\noindent \textbf{Dynamic Analysis}~\cite{nethercote2007valgrind, lehmann2019wasabi, cadar2008klee}: The goal of dynamic analysis is to evaluate the behavior of software while it is running. Running the software in a controlled environment or through automated tests while monitoring its execution and interactions with system resources is what it entails. Dynamic analysis can detect bugs in input validation~\cite{ kim2019rvfuzzer}, access control, and error handling. This approach can identify vulnerabilities that static analysis alone cannot detect by analyzing the real-time behavior of the software. However, the dynamic analysis may have constraints in terms of significant system overhead~\cite{yong2005using}.

\subsubsection{Vulnerability analysis}
After the detection of vulnerabilities, the subsequent step in software vulnerability management is vulnerability analysis and assessment~\cite{le2019automated, suciu2022expected, jacobs2023enhancing, yin2022real, le2022survey, mahor2022mobile, frei2006large}. This step involves further examining the identified vulnerabilities to assess their severity, impact, and potential exploitability.

\noindent \textbf{Severity}: Accurately assessing software vulnerabilities is vital for several reasons. Firstly, it allows organizations to prioritize their response based on the severity of the vulnerabilities. Severity refers to the potential impact a vulnerability could have if exploited~\cite{kudjo2019improving, chen2020automatic, kudjo2020effect, tan2020bug}. By accurately assessing the severity, organizations can focus their attention on high-severity vulnerabilities that pose significant risks to the security and functionality of the software system.

\noindent \textbf{Impact}: Secondly, accurately assessing vulnerabilities helps determine the potential impact they may have on the organization~\cite{gawron2018automatic, gong2019joint, chen2010categorization, jiang2020approach}. The term \textit{impact} refers to the repercussions of exploiting a vulnerability, such as denial of service~\cite{harzevili2022characterizing} or data breaches~\cite{aaltonen2021does}. By understanding the potential impact, organizations can make informed decisions regarding the urgency and priority of remediation efforts.

\noindent \textbf{exploitability}: Furthermore, accurately assessing vulnerabilities aids in understanding their potential exploitability~\cite{chen2019using, chen2019vest, bozorgi2010beyond}. This entails determining the possibility that an attacker will be successful in exploiting the vulnerability to infiltrate the software system. Organizations can estimate the amount of risk associated with each vulnerability and invest resources accordingly by evaluating criteria such as the ease of exploitation and the availability of exploit techniques.

\subsubsection{Vulnerability remediation}
The process of addressing detected software vulnerabilities by different techniques such as patching, code modification, and repairing is referred to as software vulnerability remediation~\cite{canfora2022patchworking, piantadosi2019fixing, chen2022neural, bhandari2021cvefixes}. The fundamental goal of remediation is to eliminate or mitigate vulnerabilities in order to improve the software system's security and dependability. One common approach to vulnerability remediation is applying patches provided by software vendors or open-source communities~\cite{xia2023automated, li2022generating, gissurarson2022propr}. Patches are updates or fixes that address specific vulnerabilities or weaknesses identified in a software system. 

\subsubsection{ML/DL for software vulnerability detection}

By utilizing data analysis, pattern recognition, and machine-driven learning for finding software security vulnerabilities, ML/DL approaches have revolutionized software vulnerability detection.~\cite{wang2018deep,cheng2021deepwukong,li2021vulnerability,le2021deepcva,yan2021han}. These techniques improve the accuracy and efficiency of vulnerability detection, potentially allowing for automated detection, faster analysis, and the identification of previously undisclosed vulnerabilities.

One common application of ML/DL in vulnerability detection is the classification of code snippets~\cite{dam2017automatic, shippey2019automatically, wang2016automatically, jeon2021autovas}, software binaries~\cite{phan2017convolutional,
nguyen2020deep,
yan2021han,
huang2021hunting}, or code changes mined from open-source repositories such as GitHub or CVE~\cite{sabetta2018practical,
wang2020combining,
liu2019deepbalance,
pradel2018deepbugs,
hoang2019deepjit,
zhou2019devign,
dinella2020hoppity,
3360588}. ML models can be trained on labeled datasets, where each sample represents a known vulnerability or non-vulnerability. These models then learn to generalize from the provided examples and classify new instances based on the patterns they have learned. This method allows for automatic vulnerability discovery without the need for manual examination, considerably lowering the time and effort necessary for analysis. 

ML/DL models for detecting software vulnerabilities have promising advantages over traditional methodologies. Each benefit is discussed in depth in the next paragraph. 

\noindent \textbf{Automation}: Automation is a significant advantage. ML models can automatically scan and analyze large codebases, network traffic logs, or system configurations, flagging potential vulnerabilities without requiring human intervention for each individual case~\cite{chakraborty2021deep}. This automation speeds up the detection process, allowing security teams to focus on verifying and mitigating vulnerabilities rather than manual analysis.

\noindent \textbf{Performance}: ML/DL approaches offer faster analysis. Traditional vulnerability detection methods rely on manual inspection or the application of predefined rules~\cite{staron2020using, bacchelli2013expectations, carlsson2005software, shar2012auditing, shar2010auditing}. In contrast, ML/DL approaches can evaluate enormous volumes of data in parallel and generate predictions fast, dramatically shortening the time necessary to find vulnerabilities.

\noindent \textbf{Detection effectiveness}: ML/DL models can uncover previously unknown vulnerabilities, commonly known as zero-day vulnerabilities~\cite{bilge2012before}. These models may uncover signs of vulnerabilities even when they have not been specifically trained on them by learning patterns and generalizing from labeled data. This capability improves the overall security posture by helping to identify and address unknown weaknesses in software before they are exploited by attackers~\cite{abri2019can}.




% \subsection{Definition}
% In this section, we start listing some definitions of software security vulnerabilities stated by existing studies. For example, Krsul et al.~\cite{krsul1998software} define vulnerability as \textit{An instance of an error in the specification, development, or configuration of software such that its execution can violate the security policy.} Ozment et al.~\cite{ozment2007improving} verifies Krsul's  definition of vulnerability with minor modifications: \textit{A software vulnerability is an instance of a mistake in the specification, development, or configuration of software such that its execution can violate the explicit or implicit security policy.} Ozment et al.~\cite{ozment2007improving} used the IEEE Standard Glossary of Software Engineering Terminology (IEEE Standards 1990) to justify why he has changed the word \textit{error} to \textit{mistake}. There are similar definitions provided by industry practitioners (Dowd et al.~\cite{dowd2006art}) such as: \textit{In the context of software security, vulnerabilities are specific flaws or oversights in a piece of software that allows attackers to do something malicious: expose or alter sensitive information, disrupt or destroy a system, or take control of a computer system or program.}

% One obvious pattern in the above definitions is there are multiple terms used to define what a software security vulnerability is. In order to have a general and appropriate definition of vulnerability, we need to define each term separately. In this regard, we use the IEEE Standard Glossary of Software Engineering Terminology (IEEE Standards 1990). We use the keywords \textit{error}, \textit{fault}, \textit{failure}, and \textit{mistake} and look up each one to find out the corresponding definitions. We have the following definitions:

% \begin{itemize}
%     \item \textbf{Error:} the difference between a computed, observed, or measured value or condition and the true, specified, or theoretically correct value or condition
%     \item \textbf{Fault (aka flaws or bugs):} an incorrect step, process, or data definition in a computer program.
%     \item \textbf{Failure:} the inability of a system or component to perform its required functions within specified performance requirements.
%     \item \textbf{Mistake:} a human action that produces an incorrect result.
% \end{itemize}

% Based on the above terms and their definitions, we can draw their relations as follows. When an external entity such as a human makes some mistakes, the impact is software fault (a flaw or bug) which is understandable from the difference in actual and predicted or computed results. Hence, we define a software vulnerability as \textit{A software vulnerability is a flaw or bug in a software system rooted in human mistakes in design and configuration resulting in exploitation that violates the policy and constrained defined by security researchers.}



\subsection{Related work}
%\song{For the comparison, please release more details, i.e., papers are published from year xx to year xx; how many papers are analyzed, etc}
There have been several existing survey papers on software vulnerabilities in the literature. In this section, we analyze the existing papers based on different aspects as shown in Table~\ref{relatedTable}.


\begin{table}[t!]
\caption{Comparison of contributions between our survey and the existing related surveys/reviews.}
\label{relatedTable}
\renewcommand{\arraystretch}{0.1}
\resizebox{\columnwidth}{!}{
\begin{tabular}{llccccccc}
\toprule
No &  Studies &  Data Source & Representation & Embedding & Models & Vulnerability Types & Interpretability \\
\midrule
1  & Triet et al. \cite{le2021survey}   & $$\checkmark$$ & $\times$     & $\checkmark$  & $\checkmark$ & $\times$     &  $\times$ \\
2  & Ghaffarian et al. \cite{ghaffarian2017software} & $\checkmark$ & $\checkmark$ & $\checkmark$  & $\checkmark$ & $\times$     & $\times$\\
3  & Lin et al. \cite{lin2020software} & $\checkmark$ & $\checkmark$ & $\checkmark$  & $\checkmark$ & $\times$     & $\times$\\
4  & Zeng et al. \cite{zeng2020software}  & $\checkmark$ & $\checkmark$ & $\checkmark$  & $\checkmark$ & $\times$     &  $\times$\\
5  & Semasaba et al. \cite{semasaba2020literature} & $\checkmark$ & $\checkmark$ & $\times$      & $\checkmark$ & $\checkmark$ &  $\times$\\
6  & Sun et al. \cite{sun2018data}  & $\checkmark$ & $\times$     & $\times$      & $\times$     & $\checkmark$ & $\times$\\
7  & Kritikos et al. \cite{kritikos2019survey} & $\checkmark$ & $\times$     & $\times$      & $\times$     & $\checkmark$ &  $\times$\\
8  & Khan et al. \cite{khan2018review}  & $\times$  &  $\times$  &  $\times$   &  $\times$    &  $\times$   & $\times$ \\
9  & Nong et al. \cite{nong2022open}  & $\times$  &  $\times$  &  $\times$   &  $\times$    &  $\times$   & $\times$ \\
10  & Chakraborty et al. \cite{chakraborty2021deep}  & $\times$  &  $\times$  &  $\times$   &  $\times$    &  $\times$   & $\times$ \\
11  & Liu et al. \cite{liu2021reproducibility}  & $\times$  &  $\times$  &  $\times$   &  $\times$    &  $\times$   & $\times$ \\
12  & Bi et al. \cite{bi2023benchmarking}  & $\times$  &  $\times$  &  $\times$   &  $\times$    &  $\times$   & $\times$ \\
12  & Our survey & $\checkmark$ & $\checkmark$ & $\checkmark$  & $\checkmark$ & $\checkmark$ & $\checkmark$ \\
\bottomrule
\end{tabular}}
\end{table}



The columns in the table represent different aspects of the surveys, such as the data source used, representation, feature embedding, ML/DL models, vulnerability types, and interpretability of ML/DL models. \textit{Data Source} indicates whether the survey reviewed vulnerability detection data sources. \textit{Representation} discusses whether the survey considered source code representation in its analysis. \textit{Embedding} deals with whether the survey is analyzed feature embedding in its analysis. The table also considers ML/DL models in the sixth column as \textit{ML Models}. The table also checks whether the survey considers vulnerability types based on Common Weakness Enumeration (CWE) number. The last column indicates whether the survey takes into account the interpretability of ML/DL models.%\song{You do not need to repeast your slef mutilpe time " The third " "The fourth"... "The eighth " are you serious here???}\nima{paraphrased}

%\song{you did not have "Data Source", please give reasons} \nima{one of the major significant of our work is elaboration on data sources.}

%\song{please always use "past tense" when discuss other papers, globally check all the sections!} \nima{ok}

Ghaffarian et al. \cite{ghaffarian2017software} is the closest survey to ours when it comes to data-driven security vulnerability detection. %\song{Please use one paper to explain what the check mark and what the cross mark means}\nima
{In their survey, they analyzed data-driven software vulnerability detection from various aspects including \textit{Data Sources}, \textit{Representation}, \textit{Embedding} types, and different ML/DL models as shown in Table~\ref{relatedTable}}. However, there are a couple of differences compared to our work. %First, {in terms of \textit{Vulnerability Types}, they did not cover vulnerability types in their survey, which is an important information for researchers}. Compared to Ghaffarian et al. \cite{ghaffarian2017software},
Specifically, this work also surveys vulnerability detection from the following aspects: \textbf{Comprehensive Coverage}: Understanding the many sorts of vulnerabilities allows researchers to create and develop effective vulnerability detection models that can thoroughly discover security vulnerabilities. To guarantee that their detection systems cover as many vulnerability types as possible, researchers must be familiar with the various methods of attack and potential weaknesses in software systems. \textbf{Customization of Detection Techniques}: Different sorts of vulnerabilities necessitate distinct detection methods. To build specialized detection systems that can discover certain types of vulnerabilities, researchers must first understand the subtleties of each vulnerability type. \textbf{Prioritization of Mitigation Efforts}: Researchers can prioritize mitigation efforts depending on the severity and effect of each vulnerability by understanding the many types of vulnerabilities. Critical vulnerabilities that pose the greatest danger to the system or organization can be prioritized by researchers. \textbf{Better Understanding of Attack Patterns}: Understanding the different types of vulnerabilities provides researchers with insights into the different attack patterns used by attackers. This knowledge helps researchers design detection techniques that can detect not only known attack patterns but also new, unknown patterns. 
%\nima{Second, in terms of \textit{Interpretability}, they did not cover this aspect it is important for software vulnerability detection since it allows developers and security researchers to understand how the model is making predictions and identify potential weaknesses or biases in the model's behavior}. 
\textbf{Interpretability} refers to the ability to explain how a model makes a particular decision or prediction. This is particularly important in the context of software vulnerability detection because security researchers need to be able to understand why a model is flagging a particular piece of code as potentially vulnerable. %If a DL or ML model is not interpretable, it can be difficult to identify and fix errors or biases in the model's behavior. This could result in false positives or false negatives, which could lead to missed vulnerabilities or unnecessary alerts. 
Additionally, interpretability can help improve trust in the model's predictions. If developers and security researchers can understand how a model is making its decisions, they are more likely to trust its output and take appropriate actions based on its recommendations. %Overall, interpretability is an important factor in ensuring that ML/DL models are effective tools for software vulnerability detection, and that they can be used reliably in real-world applications which is not covered in Ghaffarian et al. \cite{ghaffarian2017software} survey.

%\song{please summarize more difference, this is trivial, your paper can be directly rejected because of this!!!} \nima{done}

Triet et al.~\cite{le2021survey} reviewed data-driven vulnerability assessment and prioritization studies. They conduct a review of prior research on software assessment and prioritization that leverages machine learning and data mining methods. They examine various types of research in this area, discuss the strengths and weaknesses of each approach, and highlight some unresolved issues and potential areas for future research. The major difference to ours is that we review vulnerability detection approaches while they survey assessment and prioritization techniques. Vulnerability detection, vulnerability assessment, and vulnerability prioritization are all important components of the vulnerability management life-cycle, but they involve different stages of the vulnerability management process. Our work focuses on Vulnerability detection which refers to the process of identifying potential vulnerabilities in software systems. The goal of vulnerability detection is to identify all vulnerabilities that exist within the system, regardless of their severity. Vulnerability assessment, on the other hand, involves evaluating the severity and potential impact of each identified vulnerability. This assessment can involve analyzing factors such as the likelihood of the vulnerability being exploited and the potential harm that could result. Vulnerability prioritization involves ranking the identified vulnerabilities based on their level of risk or criticality. This ranking is typically based on the results of the vulnerability assessment, as well as other factors such as the availability of resources to address the vulnerabilities. 
%\song{cannot understand, your sentences are broken unrelated to each other} \nima{done}

Lin et al. \cite{lin2020software} examined the literature on using deep learning and neural network-based approaches for detecting software vulnerabilities. There are a couple of differences compared to our survey. First, the study of conventional source code representation techniques (Static code attributes) for software vulnerability detection. In our survey, we neglect to review such representation techniques for a couple of reasons. Static code attributes, such as code length or complexity, may not be effective for vulnerability detection because they do not capture the dynamic behavior of the code at runtime. Vulnerabilities can manifest themselves in unexpected ways that are not apparent in the static code, making it difficult to detect them through static analysis alone. Additionally, static code attributes may not be able to capture the context of the code, which is important for understanding how the code interacts with other components in a system. Finally, static analysis tools may produce a high rate of false positives, which can be time-consuming to verify and may cause developers to ignore important vulnerabilities. Second, we examine the trend analysis of papers published in software vulnerability detection in a journal and conference papers because it provides a comprehensive understanding of the publishing patterns in a particular field or area of research. The trend analysis can shed light on the distribution of research output across various publication venues and the shifting preferences of researchers and authors. This information can be useful for stakeholders such as publishers, academic institutions, and researchers in making strategic decisions related to publishing, funding, and research collaborations.

Zeng et al. \cite{zeng2020software} discussed the increasing attention towards exploitable vulnerabilities in software and the development of vulnerability detection methods, specifically the application of ML techniques. The paper reviews 22 recent studies that use deep learning to detect vulnerabilities and identifies four \textit{game changers} that have significantly impacted this field. The survey further compares the game changers based on different aspects in software vulnerability detection including data source, feature representation, DL models, and detection granularity. There are a couple of differences compared to our survey. First, we analyze the trend patterns of papers on software vulnerability detection that have been published in journals and conferences. This analysis helps us gain a thorough comprehension of the publication trends in a specific area of research or field.
Second, we cover more aspects of software vulnerability detection. While they only cover data source, feature representation, DL models, and detection granularity, we cover more aspects including vulnerability types and interpretability of ML/DL models. Additionally, we provide a more granular analysis of different aspects. 


Kritikos et al.~\cite{kritikos2019survey} and Sun et al.~\cite{sun2018data} focused on cybersecurity and aim to improve cyber resilience. Sun et al.~\cite{sun2018data} discussed the paradigm change in understanding and protecting against cyber threats from reactive detection to proactive prediction, with an emphasis on new research on cybersecurity incident prediction systems that use many types of data sources. Kritikos et al.~\cite{kritikos2019survey} discusses the challenges of migrating applications to the cloud and ensuring their security, with a focus on vulnerability management during the application lifecycle and the use of open-source tools and databases to better secure applications. While the topics of the two abstracts are different, they share a common goal of improving cybersecurity and resilience. Both highlight the importance of proactive measures to prevent or mitigate cyber threats, rather than relying solely on reactive detection and response. Additionally, both highlight the importance of utilizing various data sources and tools to improve cybersecurity measures. While both approaches aim to improve the security of applications, they differ in their focus and techniques used. They mainly focus on providing guidance and tools to support vulnerability management during the application lifecycle, while in our survey, we focus on software vulnerability detection using ML/DL techniques on source code which aim at automating the identification of vulnerabilities in the source code.

Khan et al.~\cite{khan2018review} focused on Vulnerability Assessment, which is the process of finding and fixing vulnerabilities in a computer system before they can be exploited by hackers. This highlights the necessity for more studies into automated vulnerability mitigation strategies that can effectively secure software systems. On the other hand, vulnerability identification with ML/DL approaches on source code entails analyzing an application's source code in order to spot security flaws. Instead of evaluating the safety of the entire system, this method concentrates on finding vulnerabilities in the code itself.

Nong et al.\cite{nong2022open} explored the open-science aspects of studies on software vulnerability detection and argued there is a dearth of research on problems of open science in software engineering, particularly with regard to software vulnerability detection. The authors conducted an exhaustive literature study and identify 55 relevant works that propose deep learning-based vulnerability detection approaches. They investigated open science aspects including availability, executability, reproducibility, and replicability. The study reveals that 25.5\% of the examined approaches provide open-source tools. Furthermore, some open-source tools lack adequate documentation and thorough implementation, rendering them inoperable or unreplicable. The use of unbalanced or intentionally produced datasets causes the approaches' performance to be overstated, rendering them unreplicable.


Chakraborty et al.~\cite{chakraborty2021deep} investigated the performance of cutting-edge DL-based vulnerability prediction approaches in real-world vulnerability prediction scenarios. They find that the performance of the state-of-the-art DL-based techniques drops by more than 50 percent in real-world scenarios. They also discover problems with training data (for example, data duplication and an unrealistic distribution of vulnerable classes) and model selection (for example, simplistic token-based models). Existing DL-based approaches often learn unrelated artifacts instead of features related to the cause of vulnerabilities. The significant difference compared to our survey study is that in our work, we focus on the usage of ML/DL models for software vulnerability detection and characterize the different stages in the pipeline of vulnerability detection. On the other hand, they focus on the issues related to the usage of state-of-the-art DL models for software vulnerability detection. 

Liu et al.~\cite{liu2021reproducibility} discussed the increasing popularity of DL techniques in software engineering research due to their ability to address SE challenges without extensive manual feature engineering. The authors highlight two important factors often overlooked in DL studies: reproducibility and replicability. Reproducibility refers to whether other researchers can obtain the same results using the authors' artifacts, while replicability refers to obtaining similar results with re-implemented artifacts and a different experimental setup. The major difference compared to our study is that we focus on the usage of ML/DL techniques in software vulnerability detection pipelines, while they emphasize replicability and reproducibility of the results reported in software engineering research studies. 

Bi et al.~\cite{bi2023benchmarking} emphasizes the importance of software vulnerability detection techniques as well as the absence of a systematic methodology to evaluate these approaches. The research is the first to look into and describe the current state of software vulnerability detection benchmarking. The assessment examines current literature on vulnerability detection benchmarking, including methodologies employed in technique-proposing publications and empirical research. The survey examines the difficulties associated with benchmarking software vulnerability detection approaches and suggests alternative solutions to these difficulties. They do not, however, give a characterization of datasets, representations, embedding techniques, and models employed in software vulnerability identification, unlike our work.









